% ===== PRÉAMBULE CANONIQUE — à coller VERBATIM en tête de CHAQUE .tex =====
\documentclass[11pt,a4paper]{article}
\usepackage[utf8]{inputenc}\usepackage[T1]{fontenc}\usepackage{lmodern}
\usepackage[french]{babel}
\usepackage{amsmath,amssymb,mathtools}\usepackage{icomma}
\usepackage{siunitx}\sisetup{locale=FR}  % \qty{}{}, \si{}, \ang{}
\usepackage{xcolor}\usepackage{colortbl}
\usepackage{tikz}\usepackage{pgfplots}\pgfplotsset{compat=1.18}
\usepackage[siunitx,europeanresistors]{circuitikz} % circuits (élec.)
\usepackage[most]{tcolorbox}\usepackage{enumitem}\usepackage{array,booktabs}
\usepackage[a4paper,margin=2.2cm]{geometry}
\usetikzlibrary{arrows.meta,calc,angles,quotes,positioning,patterns,%
  backgrounds,shapes.geometric,decorations.pathmorphing,decorations.markings,intersections}
\definecolor{primary}{RGB}{222,49,99}\definecolor{primarydark}{RGB}{150,22,55}
\definecolor{defcol}{RGB}{0,114,178}\definecolor{methcol}{RGB}{52,92,148}
\definecolor{excol}{RGB}{0,135,90}\definecolor{attncol}{RGB}{200,40,55}
\tcbset{boxbase/.style={enhanced,breakable,boxrule=0.9pt,arc=2.5pt,
  left=3mm,right=3mm,top=2.3mm,bottom=2.3mm,fonttitle=\bfseries,coltitle=white}}
% --- boîtes TUTORIEL (noms EXACTS, ne pas renommer) ---
\newtcolorbox{cle}[1][]{boxbase,colback=defcol!6,colframe=defcol,
  title={\textsc{À retenir}\ifx\relax#1\relax\relax\else\textmd{ : \itshape #1}\fi}}
\newtcolorbox{methode}[1][]{boxbase,colback=methcol!7,colframe=methcol,sharp corners,
  title={\textsc{Méthode}\ifx\relax#1\relax\relax\else\textmd{ --- \itshape #1}\fi}}
\newtcolorbox{exemplebox}[1][]{boxbase,colback=excol!5,colframe=excol,
  title={\textsc{Exemple}\ifx\relax#1\relax\relax\else\textmd{ : \itshape #1}\fi}}
\newtcolorbox{piege}[1][]{boxbase,colback=attncol!6,colframe=attncol,
  title={\textsc{Piège}\ifx\relax#1\relax\relax\else\textmd{ : \itshape #1}\fi}}
\newtcolorbox{remarque}[1][]{boxbase,colback=black!4,colframe=black!45,
  title={\textsc{Remarque}\ifx\relax#1\relax\relax\else\textmd{ : \itshape #1}\fi}}
% --- boîtes EXERCICE / CORRIGÉ (noms EXACTS, auto-comptées) ---
\newtcolorbox[auto counter]{exo}[1][]{boxbase,colback=primary!4,colframe=primary,
  title={\textsc{Exercice}~\thetcbcounter\ifx\relax#1\relax\relax\else\textmd{ --- \itshape #1}\fi}}
\newtcolorbox[auto counter]{corrige}[1][]{boxbase,colback=excol!4,colframe=excol,
  title={\textsc{Corrigé}~\thetcbcounter\ifx\relax#1\relax\relax\else\textmd{ --- \itshape #1}\fi}}
\newcommand{\vv}[1]{\vec{#1}}\newcommand{\ds}{\displaystyle}
\newcommand{\R}{\mathbb{R}}\newcommand{\N}{\mathbb{N}}\newcommand{\Z}{\mathbb{Z}}
\begin{document}
\sisetup{per-mode=symbol}

\begin{corrige}[vergences et pièges d'unités]
Toujours $f$ \textbf{en mètres} avant $C=1/f$ ; le signe de $C$ suit celui de $f$.
\begin{enumerate}[label=\alph*)]
  \item $f=+\qty{0.05}{\m}$ : $C=\dfrac{1}{0,05}=+20~\delta$.
  \item $f=-\qty{0.40}{\m}$ : $C=\dfrac{1}{-0,40}=-2{,}5~\delta$.
  \item $f=+\qty{2}{\mm}=+\qty{0.002}{\m}$ : $C=\dfrac{1}{0,002}=+500~\delta$ (lentille très puissante).
\end{enumerate}
\end{corrige}

\begin{corrige}[une convergente et une divergente accolées]
\begin{enumerate}[label=\alph*)]
  \item $C_{\text{éq}}=C_1+C_2=+6+(-2)=+4~\delta$.
  \item $C_{\text{éq}}>0$ : système \textbf{convergent},
        $f_{\text{éq}}=\dfrac{1}{4}=+\qty{0.25}{\m}=+\qty{25}{\cm}$.
\end{enumerate}
\end{corrige}

\begin{corrige}[fabriquer un verre neutre (afocal)]
Afocal $\Leftrightarrow C_{\text{éq}}=0$, donc $C_1+C_2=0$.
\begin{enumerate}[label=\alph*)]
  \item $C_2=-C_1=-3~\delta$.
  \item $C_2<0$ : lentille \textbf{divergente}, $f_2=\dfrac{1}{-3}=-\qty{0.33}{\m}\approx-\qty{33}{\cm}$.
        Les deux lentilles se « compensent » exactement.
\end{enumerate}
\end{corrige}

\begin{corrige}[classer par puissance]
\begin{enumerate}[label=\alph*)]
  \item $C_1=\dfrac{1}{0,20}=+5~\delta$ ; $C_2=\dfrac{1}{0,05}=+20~\delta$ ;
        $C_3=\dfrac{1}{-0,10}=-10~\delta$ ; $C_4=\dfrac{1}{1}=+1~\delta$.
  \item Convergentes ($C>0$) classées par vergence décroissante :
        \[ L_2\,(+20~\delta)\ >\ L_1\,(+5~\delta)\ >\ L_4\,(+1~\delta). \]
        $L_3$ est divergente ($C<0$). La plus puissante est $L_2$ (plus courte focale).
\end{enumerate}
\end{corrige}

\begin{corrige}[aller-retour $f \leftrightarrow C$]
$f=\dfrac{1}{C}$ en mètres, puis conversion.
\begin{enumerate}[label=\alph*)]
  \item $f=\dfrac{1}{-2,5}=-\qty{0.40}{\m}=-\qty{40}{\cm}$ : $C<0\Rightarrow$ \textbf{divergente}.
  \item $f=\dfrac{1}{20}=+\qty{0.05}{\m}=+\qty{5}{\cm}=+\qty{50}{\mm}$ : convergente et très puissante.
\end{enumerate}
\end{corrige}

\end{document}
